\section{强化学习概览}

\subsection{什么是强化学习}

\indent 典型的强化学习过程如下：

\begin{figure}[H]
  \centering
  \begin{tikzpicture}[
    x=1.16cm,
    y=1.08cm,
    >=Latex,
    font=\small,
    block/.style={
      draw=black,
      line width=0.9pt,
      rounded corners=6pt,
      align=center,
      fill=white
    },
    flow/.style={
      draw=black,
      line width=0.9pt,
      -Latex,
      line cap=round,
      line join=round
    },
    timecut/.style={draw=black, dashed, line width=0.35pt},
    label/.style={font=\small, align=center}
  ]
    \node[block, minimum width=1.95cm, minimum height=0.76cm] (agent) at (4.85,3.30) {智能体};
    \node[block, minimum width=3.40cm, minimum height=0.76cm] (env) at (4.85,0.85) {环境};

    \coordinate (stateIn) at ($(agent.west)+(0,0.16)$);
    \coordinate (rewardIn) at ($(agent.west)+(0,-0.16)$);
    \coordinate (envRewardOut) at ($(env.west)+(0,0.20)$);
    \coordinate (envStateOut) at ($(env.west)+(0,-0.20)$);
    \coordinate (timeReward) at (2.45,1.05);
    \coordinate (timeState) at (2.45,0.65);

    \draw[flow] (timeState) -- (1.00,0.65) -- (1.00,3.46) -- (stateIn);
    \node[label, anchor=east] at (0.78,2.12) {状态\\$S_t$};

    \draw[flow] (timeReward) -- (1.45,1.05) -- (1.45,3.14) -- (rewardIn);
    \node[label, anchor=west] at (1.65,2.10) {奖励\\$R_t$};

    \draw[flow] (agent.east) -- (7.55,3.30) -- (7.55,0.85) -- (env.east);
    \node[label, anchor=west] at (7.77,2.10) {动作\\$A_t$};

    \draw[flow] (envRewardOut) -- (timeReward);
    \draw[flow] (envStateOut) -- (timeState);
    \draw[timecut] (2.45,0.24) -- (2.45,1.46);
    \node[label, anchor=west] at (2.56,1.28) {$R_{t+1}$};
    \node[label, anchor=west] at (2.56,0.42) {$S_{t+1}$};
  \end{tikzpicture}
  \caption{强化学习中智能体与环境的交互过程}
  \label{fig:rl-agent-environment-loop}
\end{figure}


\begin{definitionbox}[强化学习]
A computational approach to learning whereby \textcolor{red}{an agent} tries
to \textcolor{red}{maximize} the total amount of \textcolor{red}{reward} it receives while
interacting with a complex and uncertain \textcolor{red}{environment}.

一种\textcolor{red}{智能体}在与复杂且不确定的\textcolor{red}{环境}交互时，
试图\textcolor{red}{最大化}其获得的总\textcolor{red}{奖励}的方法。

\hfill - Sutton and Barto
\end{definitionbox}

\subsection{与其他方法的区别}

\begin{enumerate}[
  label=（\arabic*）,
  labelindent=\parindent,
  labelwidth=!,
  leftmargin=*,
  labelsep=0.4em,
  parsep=0.5em
]
  \item 有监督学习（以图像分类为例）
  
  \begin{itemparagraph}
  可以认为是给出一系列标注正确答案的图片，让AI学习，
  自己找到其中的规律，然后给训练好的AI没见过的图片，让它给出答案。
  \end{itemparagraph}

  \begin{itemize}
    \item 已标注的图像（Annotated）
    \item 数据之间独立同分布（i.i.d）
  \end{itemize}
  \item 强化学习
  
  \begin{itemparagraph}
  AI进行游戏，没有引导，只能知道在游戏规则下的得分，
  AI自己学习使用什么策略方法取得尽可能高的分数。
  \end{itemparagraph}

  \begin{itemize}
    \item 非即时反馈
    \item 无正确标签
    \item 数据之间非独立同分布（相关的时间序列数据）
  \end{itemize}

  \item 对比区别
  
  \begin{itemparagraph}
  强化学习输入的是一组序列数据，并非独立同分布的数据；不存在监督者
  或者标签告诉智能体正确答案，智能体只能收到规则下的奖励信号（分数），
  并且信号也不是事实反馈的，而是存在延迟；需要通过尝试探索来发现
  给出最大奖励的策略/方法。

  强化学习中的智能体可以使用目前已知的最优策略，也需要去探索尝试
  新的策略，但尝试也有可能发现的是更差的策略，因此强化学习的一个关键点
  是\textbf{探索（exploration）与利用（exploitation）的平衡}。
  \end{itemparagraph}
\end{enumerate}

\begin{keybox}[强化学习特点]
  \begin{itemize}
    \item 序列数据
    \item 无监督学习
    \item 探索与利用的平衡
  \end{itemize}
\end{keybox}

\subsection{强化学习的优点}

\indent 当使用有监督学习的方法时，算法的上限（Upper Bound）
由人的标定的准确性所决定，不可能超过人类。而强化学习让智能体自行
在环境中进行探索，因此能够获得超人类的表现（比如Alpha GO超越了人类）。

\subsection{强化学习示例}

下面以 DeepMind 在 Atari Pong 上的经典实践为例，详细推导过程
与完整代码参见\cref{app:pong}。

\begin{figure}[H]
  \centering
  \begin{tikzpicture}[
    x=1cm,
    y=1cm,
    >=Stealth,
    font=\small,
    label/.style={font=\small},
    arrow/.style={
      draw=noteBlue,
      line width=1.1pt,
      -Stealth
    },
    pixframe/.style={
      draw=black!65,
      very thick,
      fill=brown!70!black
    },
    policybox/.style={
      draw=black,
      line width=1.0pt,
      rounded corners=4pt,
      fill=noteBlue,
      minimum width=2.65cm,
      minimum height=1.22cm,
      align=center
    }
  ]
    \begin{scope}[shift={(1.65,1.70)}]
      \node[pixframe, minimum width=3.30cm, minimum height=3.30cm] (screen) at (0,0) {};

      \draw[orange!45, line width=1.25pt]
        (-0.72,1.13) rectangle (-0.56,1.48);
      \draw[green!60!black, line width=1.25pt]
        (0.72,1.13) rectangle (0.88,1.48);

      \draw[white, line width=5.0pt]
        ([yshift=-0.76cm]screen.north west) --
        ([yshift=-0.76cm]screen.north east);

      \fill[orange!20]
        ([xshift=-1.22cm,yshift=-0.98cm]screen.center) rectangle ++(0.11,0.42);
      \fill[green!55!black]
        ([xshift=1.18cm,yshift=-1.36cm]screen.center) rectangle ++(0.11,0.42);
      \fill[white]
        ([xshift=-0.05cm,yshift=-0.18cm]screen.center) rectangle ++(0.045,0.045);

      \node[label, anchor=south] at ([yshift=7pt]screen.north) {原始像素};
    \end{scope}

    \draw[arrow] (3.55,1.70) -- (4.65,1.70);
    \node[label, anchor=south] at (4.10,1.82) {输入};

    \node[policybox] (policy) at (6.20,1.70) {\color{white}策略网络};

    \draw[arrow] (7.70,1.70) -- (8.78,1.70);
    \node[label, anchor=south] at (8.24,1.82) {输出};

    \node[label, align=left, anchor=west] at (9.00,1.88)
      {上升或下降\\的概率};
  \end{tikzpicture}
  \caption{策略网络根据 Pong 原始像素输出上升或下降的概率}
  \label{fig:pong-policy-network}
\end{figure}

